\documentclass{article}
\usepackage{amsmath}
\usepackage{amssymb}
\usepackage{fontspec}
\setlength{\parindent}{0bp}
\usepackage{xcolor}
\definecolor{lyxboxbgcolor}{rgb}{0.800781, 0.800781, 0.800781}
\usepackage{float}
\usepackage{units}
\usepackage{cancel}
\usepackage{graphicx}
\usepackage{geometry}
\geometry{verbose,tmargin=2.7cm,bmargin=2cm,lmargin=2cm,rmargin=2cm,headheight=2cm,headsep=0cm,footskip=0cm}
\usepackage{setspace}
\usepackage{microtype}
\onehalfspacing
\usepackage[pdfusetitle,
 bookmarks=true,bookmarksnumbered=true,bookmarksopen=true,bookmarksopenlevel=2,
 breaklinks=true,pdfborder={0 0 0},pdfborderstyle={},backref=section,colorlinks=true]
 {hyperref}
\hypersetup{
 linkcolor=black, citecolor=black}

\makeatletter

%%%%%%%%%%%%%%%%%%%%%%%%%%%%%% LyX specific LaTeX commands.
\providecommand\textquotedblplain{%
  \bgroup\addfontfeatures{Mapping=}\char34\egroup}

\@ifundefined{date}{}{\date{}}
%%%%%%%%%%%%%%%%%%%%%%%%%%%%%% User specified LaTeX commands.
% Engine guard: use XeLaTeX/LuaLaTeX
\usepackage{iftex}
\ifPDFTeX
  \PackageError{template}{Please compile with XeLaTeX or LuaLaTeX}{}
\fi

% ======== ADVANCED HEBREW ACADEMIC TEMPLATE ========
% Packages (trimmed to avoid conflicts)
% \usepackage{autobreak}  % REMOVE: can conflict with amsmath; use \allowdisplaybreaks
% \usepackage[T1]{fontenc}  % REMOVE: not needed with fontspec
\usepackage{relsize}
\usepackage{xcolor}
\usepackage{amssymb}
\usepackage{titlesec}
\usepackage{polyglossia}
\usepackage{fancyhdr}
\usepackage{tcolorbox}
\usepackage{tikz}
\usepackage{pgfplots}
\usepackage{booktabs}
\usepackage{array}
\usepackage{multirow}
\usepackage{enumitem}
\usepackage{float}
\usepackage{caption}
\usepackage{subcaption}
\usepackage{geometry}
\usepackage{microtype}
\usepackage{fontspec}
\usepackage{physics}
\usepackage{mathtools}
\usepackage{empheq}
\usepackage{mdframed}

% Colors (deduplicated)
\definecolor{primaryblue}{RGB}{33,150,243}
\definecolor{secondaryblue}{RGB}{63,81,181}
\definecolor{accentgreen}{RGB}{76,175,80}
\definecolor{warningcolor}{RGB}{255,152,0}
\definecolor{errorcolor}{RGB}{244,67,54}
\definecolor{infocolor}{RGB}{3,169,244}
\definecolor{notecolor}{RGB}{255,193,7}
\definecolor{theoremcolor}{RGB}{156,39,176}
\definecolor{definitioncolor}{RGB}{255,87,34}
\definecolor{examplecolor}{RGB}{0,150,136}
\definecolor{lightgray}{RGB}{245,245,245}
\definecolor{darkgray}{RGB}{66,66,66}
% Add aliases for colors referenced in layout to avoid missing-color errors
\colorlet{successcolor}{accentgreen}
\colorlet{warningred}{errorcolor}
\colorlet{purple}{theoremcolor}
\colorlet{green}{accentgreen}
\colorlet{blue}{primaryblue}
\colorlet{orange}{warningcolor}
\definecolor{brown}{RGB}{121,85,72}

% Language + fonts (polyglossia)
\setdefaultlanguage{hebrew}
\setotherlanguage{english}

% Robust font selection with fallbacks (fixes Times New Roman not loadable)
% Hebrew serif
\IfFontExistsTF{David CLM}
  {\newfontfamily\hebrewfont[Script=Hebrew]{David CLM}}
% Hebrew sans

% Math font (unchanged)
% Use Unicode math with Libertinus as default; several alternatives available.
\usepackage{unicode-math}
\setmathfont{XITS Math} % default math font
\geometry{a4paper, margin=1in}

% Setting up font for Hebrew (if using XeLaTeX or LuaLaTeX, polyglossia is preferred, 
% but babel is used here for standard compatibility)
\selectlanguage{hebrew}

\title{סיכום הרצאה: משוואת לפלס - גיאומטריות שונות}
\author{סיכום מבוסס תמליל}
\date{}

\begin{document}

\maketitle

\section{חזרה: משוואת לפלס בחצי מישור}
בשיעור הקודם עסקנו בבעיית חצי מישור ($y > 0$, $-\infty < x < \infty$) עם תנאי שפה:
$$ u(x,0) = f(x) $$
ודרישה שהפוטנציאל דועך באינסוף.
הפתרון התקבל באמצעות טרנספורם פורייה והגדרת \textbf{פונקציית גרין}. הפתרון הוא קונבולוציה בין תנאי השפה לפונקציית גרין:
\begin{equation}
    u(x,y) = \int_{-\infty}^{\infty} f(\xi) G(x,y,\xi) d\xi
\end{equation}
כאשר פונקציית גרין לחצי מישור היא:
\begin{equation}
    G(x,y) = \frac{1}{\pi} \frac{y}{y^2 + (x-\xi)^2}
\end{equation}

\subsection{משמעות פיזיקלית: כלוב פראדיי}
הוזכר כי אם תנאי השפה הוא מחזורי (למשל $f(x) = A \sin(kx)$), הפתרון דועך אקספוננציאלית בכיוון $y$:
$$ u(x,y) \sim e^{-ky} \sin(kx) $$
דבר זה מסביר את עקרון הפעולה של \textbf{כלוב פראדיי}: פוטנציאל מחזורי על השפה דועך במהירות בתוך הנפח.

\section{משוואת לפלס בפרוסה אינסופית (Infinite Strip)}
הבעיה מוגדרת בתחום: $-\infty < x < \infty$ ו- $0 < y < L$.
\begin{itemize}
    \item תנאי שפה: $u(x,0) = f(x)$ ו- $u(x,L) = 0$.
    \item דרישת התאפסות ב- $x \to \pm \infty$.
\end{itemize}

\subsection{שיטת הפתרון}
מבצעים טרנספורם פורייה לפי $x$ ומקבלים משוואה דיפרנציאלית רגילה עבור $\hat{u}(k,y)$:
\begin{equation}
    \frac{d^2 \hat{u}}{dy^2} - k^2 \hat{u} = 0
\end{equation}
הפתרון הכללי הוא סכום של אקספוננטים (או סינוס היפרבולי). לאחר הצבת תנאי השפה מתקבל הפתרון במרחב התדר:
\begin{equation}
    \hat{u}(k,y) = \hat{f}(k) \frac{\sinh(k(L-y))}{\sinh(kL)}
\end{equation}
בחזרה למרחב המקום, מקבלים פונקציית גרין חדשה המותאמת לגיאומטריה של הפרוסה.

\subsection{עקרון הסופרפוזיציה}
אם נתונים תנאי שפה לא הומוגניים בשני הצדדים ($u(x,0)=f(x)$ וגם $u(x,L)=g(x)$), מפרקים את הבעיה לשתי תת-בעיות:
\begin{enumerate}
    \item $u_1$: תנאי שפה $f(x)$ למטה ו-$0$ למעלה.
    \item $u_2$: תנאי שפה $0$ למטה ו-$g(x)$ למעלה.
\end{enumerate}
הפתרון המלא הוא הסכום: $u = u_1 + u_2$, מכוח הליניאריות של משוואת לפלס.

\section{משוואת לפלס בתחום סופי: מלבן}
גיאומטריה: $0 < x < L_x$, $0 < y < L_y$.
שיטת הפתרון: \textbf{הפרדת משתנים} $u(x,y) = X(x)Y(y)$.

נניח תנאי שפה לא הומוגני על צלע אחת בלבד ($y=0$) ו-0 בשאר הצלעות.
המשוואות המתקבלות:
\begin{align}
    X''(x) + \lambda^2 X(x) &= 0 \\
    Y''(y) - \lambda^2 Y(y) &= 0
\end{align}
בציר $x$ (בו יש תנאי שפה הומוגניים בשני הקצוות) מתקבלת בעיית שטורם-ליוביל עם ערכים עצמיים דיסקרטיים:
$$ \lambda_n = \frac{n\pi}{L_x} $$
הפתרון הכללי הוא טור:
\begin{equation}
    u(x,y) = \sum_{n=1}^{\infty} C_n \sin\left(\frac{n\pi x}{L_x}\right) \sinh\left(\frac{n\pi (L_y - y)}{L_x}\right)
\end{equation}
את המקדמים $C_n$ מוצאים על ידי הטלת תנאי השפה $f(x)$ על הפונקציות העצמיות (טור פורייה סינוס).

\section{משוואת לפלס בעיגול (קואורדינטות פולריות)}
עוברים לקואורדינטות $(r, \theta)$. הלפלסיאן:
\begin{equation}
    \nabla^2 u = \frac{\partial^2 u}{\partial r^2} + \frac{1}{r}\frac{\partial u}{\partial r} + \frac{1}{r^2}\frac{\partial^2 u}{\partial \theta^2} = 0
\end{equation}
מבצעים הפרדת משתנים: $u(r,\theta) = R(r)\Theta(\theta)$.

\subsection{המשוואה הזוויתית}
מתקבלת המשוואה $\Theta'' + \lambda^2 \Theta = 0$.
בגלל המחזוריות של המעגל ($\Theta(\theta) = \Theta(\theta + 2\pi)$), הערכים העצמיים חייבים להיות שלמים: $\lambda = n$ (כאשר $n=0,1,2,\dots$).
הפתרונות הם $\sin(n\theta)$ ו-$\cos(n\theta)$.

\subsection{המשוואה הרדיאלית (משוואת אוילר)}
\begin{equation}
    r^2 R'' + r R' - n^2 R = 0
\end{equation}
הפתרונות הם מהצורה $R(r) \sim r^\alpha$.
עבור $n \neq 0$, הפתרון הכללי הוא:
$$ R_n(r) = C_1 r^n + C_2 r^{-n} $$
עבור $n = 0$, הפתרון הוא קבוע (או $\ln r$).

\subsection{בעיה פנימית}
עבור פתרון בתוך עיגול מלא ($0 \le r < a$), הפוטנציאל חייב להיות סופי בראשית ($r=0$). לכן, זורקים את האיברים המתבדרים ($r^{-n}$ ו-$\ln r$).
הפתרון הכללי לבעיה פנימית:
\begin{equation}
    u(r,\theta) = A_0 + \sum_{n=1}^{\infty} r^n (A_n \cos(n\theta) + B_n \sin(n\theta))
\end{equation}
את המקדמים מוצאים בעזרת טור פורייה של תנאי השפה על היקף המעגל.

\end{document}
